% Part: history
% Chapter: biographies 
% Section: emmy-noether

\documentclass[../../../include/open-logic-section]{subfiles}

\begin{document}

\olfileid{his}{bio}{noe}

\olsection{એમી નોથર}

\olphoto{noether-emmy}{એમી નોથર}

એમી નોથર (ઉચ્ચાર: નર-ટર)નો જન્મ 23 માર્ચ
1882ના રોજ જર્મનીના એરલાંગનમાં ઉચ્ચ
મધ્યમવર્ગીય વિદ્વાન પરિવારમાં થયો હતો.
“આધુનિક બીજગણિતની માતા” ગણાતાં નોથરે
સ્ત્રીઓના શિક્ષણમાં મોટા અવરોધો હોવા છતાં
ગણિત અને ભૌતિકશાસ્ત્ર બંનેમાં પાયાનું
પ્રદાન કર્યું. તે સમયે જર્મનીમાં છોકરીઓને
કલાઓનું શિક્ષણ આપવાનું માનવામાં આવતું અને
તેમને કોલેજની તૈયારી કરાવતી શાળાઓમાં
પ્રવેશ મળતો ન હતો. તેમ છતાં ગોટિંગેન
અને એરલાંગન યુનિવર્સિટીઓમાં નિયમિત
પ્રવેશ વિના વ્યાખ્યાનો સાંભળ્યા પછી
(એરલાંગનમાં તેમના પિતા ગણિતના પ્રોફેસર
હતા), સ્ત્રી વિદ્યાર્થીઓને પ્રવેશની
મંજૂરી આપતી નીતિ લાગુ પડતાં તેઓ
1904માં એરલાંગનમાં વિદ્યાર્થી તરીકે
નોંધાઈ શક્યાં. 1907માં તેમણે ગણિતમાં
ડૉક્ટરેટ મેળવી.

લાયકાત હોવા છતાં નોથરને કારકિર્દીમાં
ઘણો વિરોધ સહન કરવો પડ્યો. 1908--1915
દરમિયાન તેમણે એરલાંગનમાં પગાર વગર
અધ્યાપન કર્યું. આ દરમિયાન તે સમયના
વિશ્વના અગ્રણી ગણિતશાસ્ત્રીઓમાંના
એક ડેવિડ હિલ્બર્ટનું ધ્યાન તેમના
તરફ ગયું, અને તેમણે નોથરને ગોટિંગેન
બોલાવ્યાં. પરંતુ સ્ત્રીઓને પ્રોફેસરપદ
મેળવવાની મનાઈ હતી. તેથી નોથર
હિલ્બર્ટના નામે જ, ફરી પગાર વિના,
વ્યાખ્યાનો આપી શક્યાં. આ દરમિયાન
તેમણે આજે નોથરનું પ્રમેય કહેવાતું
પરિણામ સાબિત કર્યું, જે આજે પણ
સૈદ્ધાંતિક ભૌતિકશાસ્ત્રમાં વપરાય છે.
આખરે 1919માં તેમને અધ્યાપનનો
અધિકાર મળ્યો. યુનિવર્સિટીના
સહકર્મીઓના સતત વિરોધ સામે
હિલ્બર્ટે કથિત રીતે કહ્યું હતું:
“સજ્જનો, અધ્યાપકમંડળની સેનેટ
સ્નાનગૃહ નથી.”

1920ના દાયકાના ઉત્તરાર્ધમાં તેમણે
અમૂર્ત બીજગણિત પર ધ્યાન કેન્દ્રિત
કર્યું અને તેમના પ્રદાને આ ક્ષેત્રમાં
મૂળભૂત પરિવર્તન આણ્યું. સાબિતીઓમાં
તેઓ ઘણી વાર કહેવાતી ચડતી શૃંખલાની
શરત વાપરતાં; તેનો અર્થ એવો છે કે
અમુક ગણોની અનંત, સખત રીતે ચડતી
શૃંખલા હોતી નથી. દાખલા તરીકે,
આજે નોથરિયન વલયો કહેવાતી અમુક
બીજગણિતીય સંરચનાઓમાં આદર્શોની
અનંત શ્રેણી $I_1
\subsetneq I_2 \subsetneq \dots$
હોતી નથી. આ શરતને કોઈ પણ
આંશિક ક્રમ સુધી સામાન્ય બનાવી
શકાય છે (બીજગણિતમાં તે સમાવિષ્ટ
સંબંધથી ક્રમબદ્ધ આદર્શોના ખાસ
કિસ્સા સાથે સંબંધિત છે). તેની
દ્વૈત ઊતરતી શૃંખલાની શરત પણ
વિચારી શકીએ: આંશિક ક્રમમાં દરેક
સખત રીતે \emph{ઊતરતી} શ્રેણી
છેવટે અટકી જાય. જો કોઈ આંશિક
ક્રમ ઊતરતી શૃંખલાની શરત સંતોષે,
તો $<$ ક્રમ પર~$\Nat$ માટે
જેવી રીતે અનુમાન વાપરીએ છીએ,
તેવી રીતે આ ક્રમ પર પણ અનુમાન
વાપરી શકાય. આવા ક્રમોને
\emph{સુઆધારિત} અથવા
\emph{નોથરિયન} કહે છે અને
તેના અનુરૂપ સાબિતીના સિદ્ધાંતને
\emph{નોથરિયન અનુમાન} કહે છે.

નોથર યહૂદી હતાં. 1933માં નાઝીઓ
સત્તા પર આવતાં તેમને પોતાના
પદેથી બરતરફ કરાયાં. સદનસીબે
તેમને અમેરિકાના પેન્સિલ્વેનિયામાં
બ્રિન મૉર ખાતે કામચલાઉ પદ
મળતાં તેઓ ત્યાં સ્થળાંતર કરી
શક્યાં. ત્યાં રહેતાં તેમણે
પ્રિન્સ્ટનમાં પણ વ્યાખ્યાનો આપ્યાં,
જોકે તેમને એ યુનિવર્સિટી
સ્ત્રીઓ માટે પ્રતિકૂળ લાગી
\citep[81]{Dick1981}. 1935માં
ગર્ભાશયની ગાંઠ દૂર કરવા
તેમની શસ્ત્રક્રિયા થઈ. તેના
પરિણામે થયેલા ચેપથી તેમનું
અવસાન થયું; તેમને બ્રિન
મૉરમાં દફનાવવામાં આવ્યાં.

\begin{reading} 
નોથરના જીવનચરિત્ર માટે
\citet{Dick1981} જુઓ. સૈદ્ધાંતિક
ભૌતિકશાસ્ત્રની પેરિમિટર
સંસ્થાએ તેમના જીવન અને
પ્રભાવ વિશેનાં વ્યાખ્યાનો
ઑનલાઇન મૂક્યાં છે
\citep{Perimeter2015}. વાંચીને
થાક્યા હો, તો
\emph{Stuff You Missed in History Class}
માં તેમના જીવન અને પ્રભાવ
વિશેનો પૉડકાસ્ટ છે
\citep{Frey2015}. નોથરનાં
એકત્રિત લખાણો મૂળ જર્મન
ભાષામાં ઉપલબ્ધ છે
\citep{Noether1983}.
\end{reading}

\end{document}
